\sectionbanner{一、单选题}{参考答案与解析}

\begin{solution}{1}{C}{6.0}{排列组合与概率}
\analysis{把“必然出现两种不同颜色分别达到4个和3个”反过来研究。尚未满足目标时，要么四种颜色都不满4个，要么只有一种颜色达到4个而其余三种颜色都不满3个。分别求这两类坏状态最多能取多少个球，再加1。}
\begin{enumerate}[label=\textbf{Step \arabic*:},leftmargin=*]
\item 先考虑没有任何颜色取到4个的情形。此时每种颜色至多取到3个，所以总数至多为
\[
  3+3+3+3=12.
\]
\item 再考虑已经有一种颜色至少取到4个，但仍没有另一种不同颜色达到3个的情形。达到4个的那种颜色至多有10个，其余三种颜色各至多有2个，所以总数至多为
\[
  10+2+2+2=16.
\]
这两类已经覆盖所有不满足题目目标的取法，因此坏状态最多只能取16个球。所以取17个球时，必有一种颜色至少取到4个，并且另有一种不同颜色至少取到3个。
\item 取完某一种颜色的10个球，再从其余三种颜色中各取2个，共取16个球，仍不满足要求。因此16个确实不够，最少应取17个，故C正确。
\end{enumerate}
选项辨析：A只保证第一种颜色达到4个；B把最大坏状态数误当成保证数；D在处理第二种颜色的阈值时多留了容量。
\methodsummary{复合抽屉保证题要先否定目标，把所有坏状态分完类，再用“坏状态最大值加1”得到保证数，并用边界配置证明最小性。}
\end{solution}

\begin{solution}{2}{C}{6.0}{代数与方程应用}
\analysis{转入男生、转出女生的过程中，男生人数只在转入时增加，女生人数只在转出时减少。先由原来的比例恢复男生、女生的具体人数，再用转出后女生的新占比锁定最终总人数，最后回推转入的男生数。}
\begin{enumerate}[label=\textbf{Step \arabic*:},leftmargin=*]
\item 原来男生有 \(60\times60\%=36\) 人，女生有 \(60-36=24\) 人。
\item 转出4名女生后，女生还剩 \(24-4=20\) 人。此时男生占 \(75\%\)，女生占 \(25\%\)，所以此时小组总人数为
\[
  20\div25\%=80.
\]
\item 设转入男生 \(x\) 人，则
\[
  60-4+x=80,
\]
解得 \(x=24\)。
\end{enumerate}
选项辨析：A把最终总人数误当成 \(60-4=56\)；B把转出4人误当成转入；D算出了变动后的男生总数42，却把它当成转入人数。
\methodsummary{占比随人员增减变化时，先找出只受单方向影响的不变量分段，用具体人数锁定每一段的总人数，再回推变化量。}
\end{solution}

\begin{solution}{3}{B}{6.0}{代数与方程应用}
\analysis{把两条不等式都整理成同一个差 \(a-d\) 的上下界。这个严格区间能够包含 \(a-d\)，会直接迫使右端差 \(b-c\) 为正。}
\begin{enumerate}[label=\textbf{Step \arabic*:},leftmargin=*]
\item 由 \(a+b>c+d\) 得到
\[
  a-d>c-b=-(b-c).
\]
\item 由 \(a+c<b+d\) 得到
\[
  a-d<b-c.
\]
所以
\[
  -(b-c)<a-d<b-c.
\]
\item 既然同一个数 \(a-d\) 严格位于两端之间，就必须有
\[
  -(b-c)<b-c,
\]
从而 \(b-c>0\)，即 \(b>c\)。
\end{enumerate}
选项辨析：\((a,b,c,d)=(1,4,2,1)\) 满足题设，可同时否定A、C；\((a,b,c,d)=(5,3,1,5)\) 满足题设，可否定D；B必然成立。
\methodsummary{多条线性不等式判断必然关系时，可把它们化为同一差量的上下界，再利用区间非空性锁定符号。}
\end{solution}

\begin{solution}{4}{A}{8.0}{几何与空间}
\analysis{点 \(P\) 沿折线 \(B\to A\to E\to F\to G\to C\) 匀速运动，三角形 \(PBC\) 以 \(BC\) 为底，面积完全由 \(P\) 到直线 \(BC\) 的高决定。分段观察这条高即可得到图象形状。}
设正方形边长为20，剪去的小正方形边长为10，则 \(BC=20\)。
\begin{enumerate}[label=\textbf{Step \arabic*:},leftmargin=*]
\item \(P\) 从 \(B\) 沿 \(BA\) 上升到 \(A\)：高从0线性增至20，面积从0线性增至 \( \frac12\times20\times20=200 \)。
\item \(P\) 从 \(A\) 沿 \(AE\) 走到 \(E\)：高度始终为20，面积保持200。
\item \(P\) 从 \(E\) 沿 \(EF\) 下降到 \(F\)：高度从20线性降至10，面积从200线性降至100。
\item \(P\) 从 \(F\) 沿 \(FG\) 走到 \(G\)：高度始终为10，面积保持100。
\item \(P\) 从 \(G\) 沿 \(GC\) 走到 \(C\)：高度从10线性降至0，面积从100线性降至0。
\end{enumerate}
这正是选项A的分段折线：先升、再水平、再降、再水平、最后降到0。不含端点 \(B,C\) 并不改变图象趋势。
\methodsummary{动点面积图象先把底固定，再把轨迹按“等高段”切开，每段只看高的线性变化。}
\end{solution}

\begin{solution}{5}{B}{8.0}{数论与整除逻辑}
\analysis{乘积末尾的每一个0都来自一对因数 \(2\times5\)。在这六个连续整数中，偶数比含5因子的数多得多，因此末尾0的个数完全由因数5的总个数决定。}
\begin{enumerate}[label=\textbf{Step \arabic*:},leftmargin=*]
\item 末尾连续0的个数等于
\[
  \min(\text{2的总因子数},\ \text{5的总因子数}).
\]
\item 六个数中含5因子的是20和25：
\[
  20=2^2\times5,\qquad 25=5^2.
\]
因此5的总因子数为 \(1+2=3\)。
\item 仅20与24就分别提供至少两个2因子，合计至少5个2因子，远超3个。所以0的个数由5的个数决定，为3。
\end{enumerate}
\methodsummary{末尾零个数由乘积中2与5因子数的较小者决定；连续整数乘积中2因子通常充足，只需精确清点5因子。}
\end{solution}

\begin{solution}{6}{C}{8.0}{代数与方程应用}
\analysis{三个连续整数的和恰好是中间数的3倍，积也含有中间数为因子。用中间数设元，可以在乘积与和之间直接约去公共因子。}
\begin{enumerate}[label=\textbf{Step \arabic*:},leftmargin=*]
\item 设三个连续正整数为 \(m-1,m,m+1\)，其中 \(m\ge2\)。它们的和为 \(3m\)，积为 \((m-1)m(m+1)\)。
\item 由题意
\[
  (m-1)m(m+1)=33\times3m=99m.
\]
因为 \(m\ge2\) 为正数，两边约去 \(m\)，得
\[
  m^2-1=99,\qquad m^2=100,
\]
所以 \(m=10\)（舍去负值）。
\item 三个数为9、10、11。验证：
\[
  9\times10\times11=990=33\times30.
\]
最大的一个是11。
\end{enumerate}
\methodsummary{连续整数乘积问题优先用中间元对称化，再约去公共因子把三次关系降次。}
\end{solution}

\begin{solution}{7}{A}{8.0}{几何与空间}
\analysis{不确定 \(F\) 和 \(E\) 在哪两个圆的正上方，就无法求出上底的长度。斜边 \(BF\) 和 \(CE\) 恰好穿过特定圆的圆心，可用坐标精确算出。}
\begin{enumerate}[label=\textbf{Step \arabic*:},leftmargin=*]
\item 相邻两个圆心距等于半径 \(5\,\mathrm{cm}\)。8个圆心之间有7个间隔，下底
\[
  BC=7\times5=35\,\mathrm{cm}.
\]
高即为直径 \(10\,\mathrm{cm}\)。
\item 以第1个圆心为原点，点 \(B\) 的坐标是 \((0,-5)\)。直线 \(BF\) 经过第2个圆心 \((5,0)\)，斜率为1，方程为
\[
  y=x-5.
\]
点 \(F\) 位于最高点所在的水平切线 \(y=5\) 上，代入得 \(x=10\)，即点 \(F\) 恰好位于第3个圆的正上方。由对称性，点 \(E\) 位于第6个圆的正上方。第3个圆心到第6个圆心之间有3个间隔，上底
\[
  FE=3\times5=15\,\mathrm{cm}.
\]
\item 梯形面积
\[
  \frac{(15+35)\times10}{2}=250.
\]
\end{enumerate}
\methodsummary{多个重叠规则图形中，抓住斜边穿过的圆心，用斜率把视觉交点转化为代数坐标。}
\end{solution}

\begin{solution}{8}{C}{8.0}{代数与方程应用}
\analysis{调动不改变两组的总分之和。用“乙组新总分等于原乙组总分加调出学生总分”建立第一个方程；再用调出学生平均分与留组学生平均分相差4建立第二个方程。}
\begin{enumerate}[label=\textbf{Step \arabic*:},leftmargin=*]
\item 两组总分为
\[
  24\times85=2040,\qquad 21\times79=1659.
\]
调动后乙组平均分为81，共 \(21+k\) 人，故
\[
  1659+ks=81(21+k).
\]
\item 调动后甲组余下 \(24-k\) 人，平均分为 \(s-4\)，故
\[
  2040-ks=(s-4)(24-k).
\]
\item 由第一式得 \(ks=81k+42\)，从而
\[
  s=81+\frac{42}{k}.
\]
代入第二式并整理，得
\[
  k^2-48k+252=0,\qquad (k-6)(k-42)=0.
\]
因为 \(1\le k\le24\)，故 \(k=6\)。
\item 验证：\(s=88\)，乙组新总分 \(1659+528=2187=81\times27\)；甲组余下总分 \(2040-528=1512=84\times18\)，平均分恰为 \(s-4\)。
\end{enumerate}
\methodsummary{平均数调动问题一律先转总分，求出两个根后必须用人数范围检验合法性。}
\end{solution}

\begin{solution}{9}{C}{6.0}{几何与空间}
\analysis{关键字“至少有一根”是求并集信号。把一昼夜拆成两个12小时周期，分别计算分针触发与时针触发，再用容斥扣除重叠。}
\begin{enumerate}[label=\textbf{Step \arabic*:},leftmargin=*]
\item 分针每小时在3--4之间停留5分钟（15分至20分），在8--9之间停留5分钟（40分至45分），共10分钟。12小时内分针总触发时长为 \(10\times12=120\) 分钟。
\item 12小时内，时针在3:00到4:00、8:00到9:00各触发60分钟，共120分钟。
\item 3:00--4:00期间分针两次进入触发区，重叠10分钟；8:00--9:00同理再重叠10分钟。12小时总重叠20分钟。
\item 12小时实际运作
\[
  120+120-20=220
\]
分钟。一昼夜为 \(220\times2=440\) 分钟，即7小时20分钟。
\end{enumerate}
\methodsummary{处理双指针特定区域问题，先把周期标准化，再把快变量与慢变量分层统计，最后用容斥扣除重叠。}
\end{solution}

\begin{solution}{10}{B}{8.0}{代数与方程应用}
\analysis{不逐个代入正整数 \(n\)，而是比较相邻两项 \(F(n+1)\) 与 \(F(n)\)。当 \(n\) 向右移动1时，左侧各点贡献增加、右侧各点贡献减少，二者的权重差会准确指出最低点。}
\begin{enumerate}[label=\textbf{Step \arabic*:},leftmargin=*]
\item 当 \(1\le n\le16\) 时，
\[
  F(n+1)-F(n)=n(n+1)-153.
\]
\item 当 \(n\le11\) 时，\(n(n+1)\le132<153\)，因此 \(F(n+1)<F(n)\)；当 \(n\ge12\) 时，\(n(n+1)\ge156>153\)，因此 \(F(n+1)>F(n)\)。全局唯一最低点是 \(n=12\)。
\item
\[
\begin{aligned}
F(12)
&=1\times11+2\times10+\cdots+11\times1\\
&\quad+13\times1+14\times2+\cdots+17\times5\\
&=286+235=521.
\end{aligned}
\]
\end{enumerate}
选项辨析：A是 \(F(11)=542\)；C是 \(F(13)=524\)；D是 \(F(10)=585\)。
\methodsummary{加权绝对值之和的极值可通过“向右移动1”造成的左右权重差来定位。}
\end{solution}

\sectionbanner{二、填空题}{参考答案与解析}

\begin{solution}{11}{13}{6.0}{数论与整除逻辑}
\analysis{把两个严格分数不等式转成两个正整数间隙，再用合适的倍数相加消去 \(a\)，直接得到 \(b\) 的全局下界。}
\begin{enumerate}[label=\textbf{Step \arabic*:},leftmargin=*]
\item 由 \(\dfrac{3}{8}<\dfrac{a}{b}\) 且 \(a,b\) 为正整数，得到 \(8a-3b\ge1\)；由 \(\dfrac{a}{b}<\dfrac{2}{5}\) 得到 \(2b-5a\ge1\)。
\item 将第一式乘5、第二式乘8后相加，得
\[
  b\ge5+8=13.
\]
\item 取 \((a,b)=(5,13)\) 时，
\[
  \frac{3}{8}<\frac{5}{13}<\frac{2}{5},
\]
因此下界可以达到，\(b\) 的最小值为13。
\end{enumerate}
\methodsummary{严格分数夹逼中的最小分母问题，可把两侧差值整数化，再加权消元建立下界并给出达界构造。}
\end{solution}

\begin{solution}{12}{315}{6.0}{代数与方程应用}
\analysis{乙始终单向前进，可先由乙到达270米处的时间锁定相遇时刻。再扣除甲停留的1分钟，把甲的实际行走路程写成去程加返程。}
\begin{enumerate}[label=\textbf{Step \arabic*:},leftmargin=*]
\item 乙走到270米处需要 \(270\div60=4.5\) 分钟。从甲出发起算，相遇时已经过 \(4.5+0.5=5\) 分钟。
\item 若甲仍在去程，第5分钟时应距 \(P\) 地450米；若甲正在 \(Q\) 地停留，则必须有 \(PQ=270\) 米，但甲第3.5分钟已到 \(Q\)，第4.5分钟已结束停留。因此相遇只能发生在甲返程阶段。
\item 这5分钟中甲停留了1分钟，实际行走4分钟，总路程360米。设 \(P,Q\) 相距 \(L\) 米，则
\[
  L+(L-270)=360,
\]
解得 \(L=315\)。
\item 甲到达 \(Q\) 地用时 \(3.5\) 分钟，停留到第4.5分钟，再返回0.5分钟走45米，恰到距 \(P\) 为270米处。
\end{enumerate}
\methodsummary{折返相遇题可先用运动过程较简单的一方锁定绝对时刻，再把另一方的分段路程合并成总路程方程。}
\end{solution}

\begin{solution}{13}{28}{6.0}{几何与空间}
\analysis{不必求出每一个小长方形的具体长和宽。把中间长方形 \(EFGH\) 的四条边看成被大长方形边界“挤压”后又被涂色部分“补回”的线段，整体代入即可。}
\begin{enumerate}[label=\textbf{Step \arabic*:},leftmargin=*]
\item 大长方形周长为 \(2\times(12+14)=52\)。4个角上小长方形每个周长18，故各自“长+宽”为9，总和为36。4个涂色小长方形周长和为24，故各自“长+宽”总和为12。
\item 水平方向上，中间长方形的上、下边等于大长方形的宽减去左右空白、再加上涂色宽度；竖直方向同理。
\item 四边相加后得到
\[
  EFGH\text{周长}
  =ABCD\text{周长}
  -[\text{四角长+宽总和}]
  +[\text{涂色长+宽总和}]
  =52-36+12=28.
\]
\end{enumerate}
\methodsummary{嵌套矩形周长题应把水平线段与竖直线段分组，再用整体代数替换消元。}
\end{solution}

\begin{solution}{14}{11}{8.0}{排列组合与概率}
\analysis{10个互不相同的正差都不超过 \(M\)，先得到 \(M\ge10\)。边界值 \(M=10\) 时，10个差只能恰好是 \(1\) 到 \(10\)；但五个点的全部正差之和必为偶数，与 \(55\) 冲突。}
\begin{enumerate}[label=\textbf{Step \arabic*:},leftmargin=*]
\item 因为10个正差互不相同且每个差都不超过 \(M\)，所以 \(M\ge10\)。
\item 若 \(M=10\)，则10个差只能恰好是 \(1,2,\ldots,10\)，总和为55。另一方面，把所有正差求和，得到的是所选五点的一个偶线性组合，必为偶数，不可能等于55。因此 \(M\ge11\)。
\item 当 \(M=11\) 时，选择 \(0,1,4,9,11\)，所得10个正差为
\[
  1,2,3,4,5,7,8,9,10,11,
\]
互不相同。因此最小值为 \(M=11\)。
\end{enumerate}
\methodsummary{差集极值题可先用差值容量给粗下界，再在边界值处利用全部差值之和的结构性奇偶阻断，最后用显式构造闭环。}
\end{solution}

\begin{solution}{15}{171}{8.0}{数论与整除逻辑}
\analysis{相邻三个数中，中间数要么是局部极大（LCM），要么是局部极小（GCD）。不能有两个相邻的极大值或极小值，所以12个数必定严格按照“大、小、大、小……”交替排列。}
设最顶端的5为第1个位置 \(c_1\)，顺时针依次到 \(c_{12}\)。已知5是65和55的最大公因数，故5是极小值，65和55都是极大值。
\begin{enumerate}[label=\textbf{Step \arabic*:},leftmargin=*]
\item \(c_{12}=65\) 是 \(c_{11}\) 和5的最小公倍数，故 \(c_{11}=13\)。同理 \(c_3=11\)。
\item \(c_6=56\) 是 \(c_5\) 和 \(c_7\) 的最小公倍数，且 \(c_5\le9\)。可能值为1、2、4、7、8。若 \(c_5\) 为1、2或4，则 \(c_7\) 只能是56，与 \(c_6\) 重复。因此 \(c_5\) 只能是7或8。
\item 当 \(c_5=7\) 时，\(c_7=8\)，于是 \(x=\mathrm{lcm}(8,3)=24\)。
\item 当 \(c_5=8\) 时，\(c_7\) 可以是7、14或28，于是 \(x\) 分别为21、42、84。
\item \(x\) 的所有可能值为24、21、42、84，总和
\[
  24+21+42+84=171.
\]
\end{enumerate}
\methodsummary{处理环形数论问题时，先按极大/极小交替给位置染色，再用质因数作为突破口逐格推导。}
\end{solution}

\sectionbanner{三、多选题}{参考答案与解析}

\begin{solution}{16}{A、B、C}{6.0}{排列组合与概率}
\analysis{条件只限制第一次取出的2个球，因此应把基本结果写成“第一次的无序球对＋第3个球”。先锁定15个合法首轮球对，再让每个球对连接7个等可能的第3球结果。}
\begin{enumerate}[label=\textbf{Step \arabic*:},leftmargin=*]
\item 合法球对：恰有1个黄球 \(2\times7=14\) 对；恰有2个黄球1对。共15个等可能首轮球对，其中14对恰有1个黄球，因此A正确。
\item 每个首轮球对确定后袋中都剩7个球，条件下共有 \(15\times7=105\) 个等可能结果。
\item 若首轮恰有1个黄球，则袋中还剩1个黄球，这样的首轮球对有14对；若首轮取到2个黄球，则袋中不剩黄球。所以第3个球为黄球的结果有14个，条件概率为 \(\dfrac{14}{105}=\dfrac{2}{15}\)，B正确。
\item 三球恰有1个黄球时，首轮必须恰有1个黄球，并且第3个球不是黄球，有 \(14\times6=84\) 个结果，条件概率为 \(\dfrac{84}{105}=\dfrac45\)，C正确。
\item 若3个球都不是红球，合法首轮球对有9对，每一对之后还有4个非红球可作第3个球，共36个结果，条件概率为 \(\dfrac{36}{105}\ne\dfrac13\)，D错误。
\end{enumerate}
\methodsummary{分阶段条件概率应保留阶段标签，把条件发生前后的结果写成等可能状态。}
\end{solution}

\begin{solution}{17}{A、B、C}{8.0}{排列组合与概率}
\analysis{先任取两个不同小组，它们必须恰有一名共同学生。其余所有小组都含这名学生，就形成“共同中心”结构；若出现一个不含共同学生的小组，它会被强迫成一个三人三组的三角结构。}
任意符合要求的清单，要么所有小组都含同一名学生，要么包含在某个三人三组的三角结构中。共同中心结构至多有6个小组，三角结构至多有3个小组，所以A正确。若清单至少有4个小组，就不可能是三角结构，B正确。含6个小组时必须选定一名共同学生并列出他与其余6名学生组成的全部小组，共7种，C正确。不能再加入任何新小组的清单有7个完整共同中心清单，再加上 \(\binom73=35\) 个三角清单，共42种，不是35种，D错误。
\methodsummary{两人小组两两相交时，固定两组后第三组是否包含原共同学生，会把结构分成共同中心与三角形两类。}
\end{solution}

\begin{solution}{18}{A、B、C}{8.0}{新定义与抽象建模}
\analysis{两种正向操作都保持两个坐标的最大公因数不变，所以所有可达数对的最大公因数都是1。可达问题应逆向转为减法。}
除起点外，可达数对的两个坐标不可能相等，否则最大公因数等于该坐标，只能是 \((1,1)\)，故C正确。对任意非起点可达数对，较大坐标减去较小坐标便得到唯一前驱，故B正确。对 \((8,19)\) 逆推：
\[
  (8,19)\to(8,11)\to(8,3)\to(5,3)\to(2,3)\to(2,1)\to(1,1),
\]
共6次逆操作，A正确。对 \((7,23)\) 逆推共7次，D错误。
\methodsummary{加法型可达问题常应逆向转为减法；确定的逆过程能同时解决可达性、前驱确定性和步数。}
\end{solution}

\begin{solution}{19}{A、B、C}{8.0}{排列组合与概率}
\analysis{天平一次只有三种结果，所以第一次称量必须把9个可能的重硬币分成三个不超过3个的候选组；这会反推出两盘各放3枚。}
设第一次左右两盘各放 \(k\) 枚。第二次只有三种结果，因此任何一种第一次结果下候选硬币数都不能超过3，于是
\[
  k\le3\quad\text{且}\quad 9-2k\le3,
\]
故 \(k=3\)，A正确。第一次3枚对3枚后，若平衡则在未称的3枚中一对一称量，若不平衡则在较重盘的3枚中一对一称量，三种结果恰好对应三枚，B正确。固定方案：第一次始终称 \(\{1,2,3\}\) 对 \(\{4,5,6\}\)，第二次始终称 \(\{1,4,7\}\) 对 \(\{2,5,8\}\)，两次结果的组合能唯一确定9枚中的重硬币，C正确。若有10枚硬币，第一次3对3且平衡后有4个候选，第二次无法分开，D错误。
\methodsummary{先用“一个称量只有三种结果”反推候选组容量，再用自适应或固定的三路编码完成构造。}
\end{solution}

\sectionbanner{四、综合题}{参考答案与解析}

\begin{solution}{20（1）}{5}{6.0}{新定义与抽象建模}
\analysis{以数字1为原点，向右为 \(x\) 轴正方向，向上为 \(y\) 轴正方向。奇数平方位于右上角，偶数平方加1位于左下角。}
25恰好是5的平方，刚好是一个 \(5\times5\) 正方形矩阵的右上角顶点。它所在的最右侧竖列贯穿5行，因此包含5张卡片。
\methodsummary{螺旋方阵用“平方数定位法+坐标平移”，避免盲目画图。}
\end{solution}

\begin{solution}{20（2）}{144}{6.0}{几何与空间}
\analysis{100是10的平方。偶数平方的下一格是该圈左下角。}
第5圈左下角顶点是 \(10^2+1=101\)，坐标为 \((-5,-5)\)。100在101的正右方，坐标为 \((-4,-5)\)。左下角“好”字位置为 \((-5,-6)\)。第6圈左下角顶点是 \(12^2+1=145\)，坐标为 \((-6,-6)\)，其右边一格就是144。
\methodsummary{先用平方数锁定圈层顶点，再沿该圈边缘平移一格。}
\end{solution}

\begin{solution}{20（3）}{13}{8.0}{新定义与抽象建模}
\analysis{先找最接近300的奇数平方，再沿螺旋走完剩余步数，最后在该行上找最靠近中心的数。}
\(17^2=289\) 位于第8圈右上角，坐标 \((8,8)\)。290在其右边 \((9,8)\)，再向下10步到达300，坐标 \((9,-2)\)。\(y=-2\) 这一行能碰到的最内部圈层是第2圈底部，从右向左依次为13、14、15、16、17，最小数字是13。
\methodsummary{同一横排中数字最小者必在该行所能穿过的最内圈边缘上。}
\end{solution}

\begin{solution}{21（1）}{24}{6.0}{新定义与抽象建模}
\analysis{第 \(n\) 组的分母是 \(n+1\)，组内分数个数等于 \(\varphi(n+1)\)。}
第71组的分母为72。因为 \(72=2^3\times3^2\)，
\[
  \varphi(72)=72\times\frac12\times\frac23=24.
\]
\methodsummary{第 \(n\) 组分数个数就是欧拉函数 \(\varphi(n+1)\)。}
\end{solution}

\begin{solution}{21（2）}{279}{6.0}{新定义与抽象建模}
\analysis{到第420组为止，分母从2到421。分子为3的最简真分数要求分母 \(\ge4\) 且不是3的倍数。}
区间4到421共有418个数。其中3的倍数有139个（1到420中140个3的倍数，去掉3本身）。因此
\[
  418-139=279.
\]
\methodsummary{限定分子的逆向计数，把互质条件翻译成分母不能被分子整除，再用总数减去坏情况。}
\end{solution}

\begin{solution}{21（3）}{36}{8.0}{新定义与抽象建模}
\analysis{从第20组到第30组，分母从21到31。只统计这些组中合法分子所含质因数5的个数。注意25贡献两个5；当分母本身是5的倍数时该组贡献为0。}
逐组分母21到31筛查后，各组贡献的5的个数依次为
\[
  3,2,4,1,0,4,5,4,6,0,7,
\]
总和为36。
\methodsummary{质因数计数时只枚举含该质因数的分子，并对每个分子做互质性审查。}
\end{solution}

\begin{solution}{22（1）}{21}{6.0}{新定义与抽象建模}
\analysis{直径比为 \(20:30:40=2:3:4\)，转过相同弧长时格数比为 \(6:4:3\)。}
\(A\) 顺时针旋转90度即转过6格，因而 \(C\) 转过3格。\(C\) 顺时针旋转时顶部数字变小，由24变为21。
\methodsummary{传动格数等于主动轮格数乘以直径比的倒数，再结合顺逆时针决定顶部数字加减。}
\end{solution}

\begin{solution}{22（2）}{15}{6.0}{新定义与抽象建模}
\analysis{用边缘滚过的总弧长作为统一标尺。颜料先走完 \(B\) 的半周才到达 \(C\)，再走完 \(C\) 的一周才能把 \(C\) 涂满。}
总弧长为
\[
  15\pi+40\pi=55\pi.
\]
\(C\) 周长 \(40\pi\)，故转过 \(\dfrac{11}{8}\) 圈，即33格。去掉完整的24格后还剩9格，顶部数字 \(24-9=15\)。
\methodsummary{把复杂传动全部转化为接触点滚过的总长度：延迟弧长加上目标轮周长。}
\end{solution}

\begin{solution}{22（3）}{3}{8.0}{新定义与抽象建模}
\analysis{\(A\) 只有半圈有颜色，会交替传出长度为 \(10\pi\) 的颜料带和空白带。\(C\) 第一圈会留下唯一一段空白区，必须等这段空白再次回到接触点并迎上已满涂的 \(B\)。}
以弧长 \(d\) 为时间轴：当 \(d=50\pi\) 时 \(B\) 已全部涂满。\(C\) 第一圈中唯一空白对应 \(d\in[25\pi,35\pi]\)，它在 \(d\in[65\pi,75\pi]\) 再次到达接触点。此时 \(B\) 已是全涂滚筒，空白在 \(d=75\pi\) 时被补完。\(C\) 转过 \(\dfrac{15}{8}\) 圈即45格，相当于多转21格，顶部数字 \(24-21=3\)。
\methodsummary{局部周期涂色问题应建立距离轴，分别写出源头发射区间、目标接收区间和空白重现区间，三轴比对求交。}
\end{solution}
